\section{Introduction}
\label{sec:intro}

An agentic evaluation reports a number for a model. That number is produced by an apparatus:
a transport that carries tool calls between the model and the harness, an endpoint that serves
the weights, a set of tools with a particular return format, and a decoding configuration. The
model's name appears in the table. The apparatus does not.

This paper measures how much of that number the apparatus determines rather than the model, and finds
two answers of different kinds. The transport moves a number, which is a
bias. The endpoint deletes a row, which is a selection. A bias can be bounded, and a reader who knows
its size can correct for it. A selection cannot, because the model that was
refused has no line in the table to correct.

\subsection{The confound this study removes}

A companion study~\citep{companion2026baseline} measured a 10.4pp drop when one model moved from native
function calling to a textual protocol. It is a companion manuscript
under double-anonymous review, so a reader cannot obtain it. Both figures this paper draws from it
recompute from a snapshot of the exact cells they need, released in the deposit with their per-turn logs. On a single infrastructure that number admits two readings no
further collection on that infrastructure can separate --- a property of the model, or of that
provider's shim --- because model and infrastructure are perfectly confounded when there is only one
endpoint.

Serving the same nominal model from two managed clouds tests something a single endpoint cannot: whether
the measured outcome changes across deployed services at fixed nominal model identity. It does not, by
itself, say which component of a service produced any difference it finds --- the two stacks differ in
several observable quantities at once (\S\ref{sec:runtime}) --- but a difference that survives the change
of service is a difference the model's name alone does not describe. That cell --- same nominal model, two
commercial endpoints, transport manipulated as a paired factor, on a multi-turn downstream task --- is what
this study contributes as a design, and to our knowledge it has not been run.

\subsection{Contributions}

The order below is the one in which we would have a reader weigh them, not the one in which they were
planned. \textbf{The first is the primary contribution: it does not rest on a $p$-value, and it is
checkable by re-issuing the probes we release.}

\begin{enumerate}
\item \textbf{Eight mechanisms by which a serving endpoint removes a model from an evaluation before it
  produces a score, or destroys its measurements from inside one}, each with the endpoint's message
  (\S\ref{sec:census}). They fall into four operational classes --- lifecycle and access policy,
  capability and API-contract constraint, catalogue and quota inconsistency, integration and adapter
  incompatibility --- and they sit at three levels of detection difficulty, which is the operational
  finding: four are returned to an ordinary single-request capability check, three need a deliberately
  constructed conversational state, and one appears only when the model itself emits the output that
  triggers it. For the three that a single-request check does not reach we give the constructed probe that
  does, and for the last we give the trajectories in which it fired.
\item \textbf{Temperature zero is not determinism}, measured across the corpus: one model returns
  identical runs on 91--100\% of binaries and another on 20--24\% (\S\ref{sec:results}). That output
  distributions can differ between services at fixed model identity is already
  established~\citep{gao2025equality}; in one of eight comparisons here the difference reaches the
  \emph{stability of a downstream agentic score}, which is the extension this study contributes on that
  axis.
\item \textbf{A pre-registered $4\times2\times2$ design} --- four models, two named managed clouds, two
  transports, 45 held-out binaries, 8 runs each --- with the analysis script frozen by hash
  \emph{together with} the hypotheses before collection began; two of the ten tests were implemented
  under a timestamped amendment after it had (\S\ref{sec:method}). No test passes its Holm threshold.
  We report per contrast the effect its own observed dispersion would have let it resolve, because those
  dispersions span a factor of 23.7 and no single band describes them.
\item \textbf{A variance calibration that inverts on its own data.} This design was sized expecting
  run-to-run noise to dominate; measured here, the lever that binds is the number of tasks --- 0.5\% noise
  and a floor of 668 binaries on the contrast with the largest observed dispersion
  (\S\ref{sec:rq6}). Whether such
  calibrations transfer between apparatuses is a question this raises and does not settle, and the stronger
  reading is a proposition rather than a result (\S\ref{sec:proposizioni}).
\item \textbf{An artifact} carrying every raw row including the invalidated batches, the per-turn session
  traces, the frozen documents with their hashes, and the amendment log (\S\ref{sec:artifact}).
\end{enumerate}

\subsection{The boundary of the claim}

Three boundaries fix what follows. The transport contrast measures that the choice moves the number, not
which choice is right. The eight mechanisms are existence proofs, each with its own denominator rather than
a rate, and the class is open by construction. And the effect sizes belong to the scope
\S\ref{sec:threats} bounds.

\textbf{The transferable output of this study is the checklist of Table~\ref{tab:checklist} and the
variance-decomposition method, not the coefficients.} Two of the four models are open-weight, so the task,
the corpus and the harness can be re-run on infrastructure a reader controls --- though not this design:
every cell here was served by one of the two commercial clouds, and a self-hosted re-run would be a
different apparatus, which is precisely the property under study. The propositions of
\S\ref{sec:proposizioni} are hypotheses this study instantiates, not claims it establishes. The narrower
claim made here is that an evaluation which does not report its transport is not reproducible from its
method section, whatever the task.